\documentclass[12pt]{scrartcl}
\usepackage{fontawesome5}
\usepackage{siunitx}

%%%%%%%%%%%%%%%%%%%%%%%%
%        Minted        %
%%%%%%%%%%%%%%%%%%%%%%%%

\usepackage{minted}
\setminted[python]{
    linenos,
}
\usemintedstyle{borland}

\setlength{\parindent}{0pt}
\setlength{\parskip}{1em}

%%%%%%%%%%%%%%%%%%%%%%%%%%%%%%%%%%
%        Graphics & boxes        %
%%%%%%%%%%%%%%%%%%%%%%%%%%%%%%%%%%

\usepackage{xcolor}
\usepackage{hyperref}
\usepackage{catppuccinpalette}
\usepackage{tikz}
\usepackage[breakable]{tcolorbox}
\tcbuselibrary{theorems, skins}

\tikzset{
	second box/.style={
			anchor=east,
			text=white,
			% rounded corners,
			fill=#1,
			xshift=-4mm,
		},
	tcbreak/.style = {fill=white, thick, draw=#1, text=#1, font=\bfseries}
}

\tcbset{
	common/.style n args={2}{
            parbox=false,
			colframe={#1},
			colback={#1!5},
			colbacktitle={#1},
			lower separated=false,
			coltitle=white,
			boxrule=1.25pt,
			fonttitle=\bfseries,
			enhanced,
			breakable,
			top=10pt,
			sharp corners,
			before skip=1em,
			after skip=2em,
			attach boxed title to top left={
					yshift=-0.25cm,
					xshift=0.38cm,
				},
			boxed title style={
					boxrule=0pt,
					colframe=white,
					arc=0pt,
					outer arc=0pt,
				},
			separator sign={~~},
			% -- overlays based on the following TeX SE answer: https://tex.stackexchange.com/a/545324/162087
			overlay unbroken={
					\node[text=white, align=right, fill=#1, xshift=-4mm, minimum height=6mm, anchor=east] at (frame.south east) {#2};
				},
			overlay first={
					\draw[thick, #1, decoration={coil, amplitude=0.5mm}, decorate]
					(frame.south west) -- (frame.south east);
					\path[font=\small\itshape] (frame.south) node [tcbreak={#1}] (cont) {continues in the next page \faHandPointRight};
				},
			overlay middle={%
					% upper break
					\draw[thick, #1, decoration={coil, amplitude=0.5mm}, decorate]
					(frame.north west) -- (frame.north east);
					\path[font=\small\itshape] (frame.north) node [tcbreak={#1}] (cont) {\faHandPointRight continues from the previous page};

					% lower break
					\draw[thick, #1, decoration={coil, amplitude=0.5mm}, decorate]
					(frame.south west) -- (frame.south east);
					\path[font=\small\itshape] (frame.south) node [tcbreak={#1}] (cont) {continues in the next page \faHandPointRight};
				},
			overlay last={
					\node[text=white, align=right, fill=#1, xshift=-4mm, minimum height=6mm, anchor=east] at (frame.south east) {#2};
					\draw[thick, #1, decoration={coil, amplitude=0.5mm}, decorate]
					(frame.north west) -- (frame.north east);
					\path[font=\small\itshape] (frame.north) node [tcbreak={#1}] (cont) {\faHandPointRight continues from the previous page};
				}
		},
	examplestyle/.style={
			common={CtpLatteBlue}{\faStar},
		},
}

\newcommand{\newterm}[1]{\textcolor{CtpLatteMauve}{\textbf{#1}}}

% args are {environment name}{Title}{style}{reference prefix}
\newtcbtheorem[auto counter, number within=section]{example}{Example}{examplestyle}{example}

%%%%%%%%%%%%%%%%%%%%%%%%%%%%%%%
%        Maths related        %
%%%%%%%%%%%%%%%%%%%%%%%%%%%%%%%

\usepackage{amsmath, amsfonts, amssymb}
\newcommand{\mtrx}[1]{\mathbf{#1}}
\newcommand{\vvec}[1]{\vec{\mathrm{#1}}}
\newcommand{\uvec}[1]{\hat{\mathrm{#1}}}
\newcommand{\T}[1]{#1^{\intercal}}

\colorlet{xcomp}{CtpLatteRed}
\colorlet{ycomp}{CtpLatteBlue}
\colorlet{zcomp}{CtpLatteGreen}

% Row- and column-vectors: arguments separated by ";".
% Example: $\vec{a} = \colvec{1;2;3;4}$.
\makeatletter
\newcommand\rcvector[2][\\]{\ensuremath{%
		\global\def\rc@delim{#1}%
		\negthinspace\begin{bmatrix}
			\rc@vector #2;\relax\noexpand\@eolst%
		\end{bmatrix}}}
\def\rc@vector #1;#2\@eolst{%
	\ifx\relax#2\relax
        #1
	\else
        #1\rc@delim
		\rc@vector #2\@eolst%
	\fi}
\makeatother
\newcommand{\colvecxy}[2]{%
    \begin{bmatrix}
        \textcolor{xcomp}{#1}\\
        \textcolor{ycomp}{#2}
    \end{bmatrix}
}
\newcommand{\rowvec}[1]{\rcvector[\ \;]{#1}}


%%%%%%%%%%%%%%%%%%%%%%%%%%
%        Doc data        %
%%%%%%%%%%%%%%%%%%%%%%%%%%

\title{Course Name: Exercise 1}
\subtitle{Power method for finding dominant eigenvector of a matrix}
\author{Peleg Bar Sapir}

%%%%%%%%%%%%%%%%%%%%%%%%%%%%%%%
%        Document here        %
%%%%%%%%%%%%%%%%%%%%%%%%%%%%%%%

\begin{document}
\maketitle

\section{Theoretical overview}
\subsection{Extracting the eigenvalue of a known eigenvector}
Let $\vvec{v}\in\mathbf{R}^{n}$ be an eigenvector of the matrix $\mtrx{A}$ with the corresponding eigenvalue $\lambda$. Then by normalizing $\vvec{v}$ and multiplying it with $\mtrx{A}$ from both the right and the left, we can extract $\lambda$:
\begin{equation}
    \begin{aligned}
        \T{\uvec{v}}\mtrx{A}\uvec{v} &= \T{\uvec{v}}\lambda\uvec{v}\\
                                     &= \lambda\T{\uvec{v}}\uvec{v}\\
                                     &= \lambda\left\|\uvec{v}\right\|^{2}\\
                                     &= \lambda\cdot1\\
                                     &= \lambda.
    \end{aligned}
    \label{eq:eigenval_from_vec}
\end{equation}

\subsection{The power method}
The \textit{dominant eigenvalue} $\lambda_{\mathrm{dom}}$ of a matrix $\mtrx{A}$ of dimensions $n\times n$ is the eigenvalue with largest (absolute) value:
\begin{equation}
    \lambda_{\mathrm{dom}} = \max\left|\lambda_{i}\right|,
    \label{eq:}
\end{equation}
where $\left\{\lambda_{i}\right\}$ are the eigenvalues of $\mtrx{A}$.

One can find the dominant eigenvector of $\mtrx{A}$ by using the \textit{power method}:
\begin{enumerate}
    \item let $\vvec{v}_{0}\in\mathbb{R}^{n}$ be a \textbf{non-zero} vector with random components.
    \item Set $\lambda_{0}=1$ (note: the subscript denotes an iteration step, not the order of the eigenvalue/vector).
    \item Apply $\mtrx{A}$ to $\vvec{v}_{0}$ and normalize the result: $\vvec{v}_{1}=\mtrx{A}\vvec{v}_{0},\ \uvec{v}_{1}=\frac{1}{\left\|\vvec{v}_{1}\right\|}\vvec{v}_{1}$.
    \item Extract $\lambda_{1}$ by multiplying $\mtrx{A}$ from the left by $\T{\uvec{v}_{1}}$ and from the right by $\uvec{v}_{1}$ (\autoref{eq:eigenval_from_vec}).
    \item Repeat:
        \begin{enumerate}
            \item $\vvec{v}_{i+1} = \mtrx{A}\uvec{v}_{i}$,
            \item $\uvec{v}_{i+1} = \frac{1}{\left\|\vvec{v}_{i+1}\right\|}\vvec{v}_{i+1}$,
            \item $\lambda_{i+1} = \T{\uvec{v}_{i+1}}\mtrx{A}\uvec{v}_{i+1}$,
        \end{enumerate}
        until $\left|\lambda_{i}-\lambda_{i+1}\right| < \varepsilon$ (where $0<\varepsilon\in\mathbb{R}$ is a pre-defined threshold).
\end{enumerate}

\section{Actual exercise}
Implement the power method algorithm in Python, using NumPy:
\begin{enumerate}
    \item Generate a random $5\times5$ matrix $\mtrx{A}$.
    \item Generate a random starting vector $\vvec{v}\neq\vvec{0}$.
    \item Implement the algorithm using the threshold $\varepsilon=\num{1.0e-7}$ and a maximum number of iterations $k_{\max}=\num{1000000}$ (one million).
    \item Compare your results to the standard NumPy function \mintinline{python}{numpy.linalg.eig()}.
\end{enumerate}

\textbf{Notes}:
\begin{enumerate}
    \item NumPy's \mintinline{python}{numpy.linalg.eig()} returns two variables: the eigenvalues of the input matrix as a single vector, and the corresponding eigenvectors as columns of a single matrix.
    \item You can use the notation \mintinline{bash}{A@B} to calculate the product of two matrices \mintinline{python}{A} and \mintinline{python}{B}.
    \item You can use \mintinline{python}{A.T} to get the transpose of the matrix \mintinline{python}{A}.
\end{enumerate}

\end{document}
